% year: תשפ"ג
% type: מבחן
% semester: א
% moed: א
% number: 4ב
% points: 10
% categories: taylor
% source: exams/מבחנים/תשפג/מועד א תשפג.pdf

\begin{question}
(10 נק') חשבו את $\sin 0.1$ בשגיאה שלא תעלה על $0.001$.
\end{question}

\begin{hint}
השתמשו בפולינום מקלורן של $\sin x$ והעריכו את השארית בצורת לגרנז'.
\end{hint}

\begin{solution}
נשתמש בפולינום מקלורן של $f(x)=\sin x$. מתקיים $f(0)=0$, $f'(0)=\cos 0=1$, $f''(0)=-\sin 0=0$, ולכן
\[ P_2(x)=x. \]
לפי נוסחת טיילור עם שארית לגרנז', קיימת $c$ בין $0$ ל-$x$ כך ש-
\[ \sin x=P_2(x)+R_2(x),\qquad R_2(x)=\frac{f'''(c)}{3!}x^3=\frac{-\cos c}{6}x^3. \]
עבור $x=0.1$, ומכיוון ש-$|\cos c|\le 1$:
\[ |R_2(0.1)|\le\frac{(0.1)^3}{6}=\frac{0.001}{6}\approx 0.000167<0.001. \]
לכן $\sin 0.1\approx 0.1$, בשגיאה קטנה מ-$0.001$ כנדרש.

(לקירוב מדויק יותר אפשר לקחת $P_4(x)=x-\frac{x^3}{6}$, ואז $\sin 0.1\approx 0.0998333$ בשגיאה $\le\frac{(0.1)^5}{120}<10^{-7}$; הערך האמיתי $0.0998334\ldots$.)
\end{solution}
